\documentclass{article}
\usepackage[margin=1.5cm]{geometry}
\usepackage{amsmath}
\usepackage{enumitem}

\setlist{nosep,leftmargin=*}

\begin{document}
\twocolumn

\title{Chapter 6 In-Class Exercise: Building a Fourier-Series Program}
\author{Prof. Jordan C. Hanson}
\date{}

\maketitle

\section{Step 1: One Fourier Term}

\begin{enumerate}

\item The sawtooth approximation is

\[
f_N(x)=
2\sum_{n=1}^{N}
\frac{(-1)^{n+1}}{n}\sin(nx).
\]

Begin by writing a function that evaluates only the
$n$-th term:

\begin{verbatim}
import numpy as np
import matplotlib.pyplot as plt

def fourier_term(x, n):
    term = 2 * (-1)**(n + 1)
    term = term * np.sin(n*x) / n
    return term
\end{verbatim}

Test the function:

\begin{verbatim}
x = np.pi / 4

print(fourier_term(x, 1))
print(fourier_term(x, 2))
print(fourier_term(x, 3))
\end{verbatim}

\begin{itemize}
\item What are the two parameters of the function?
\item What changes when \texttt{n} changes?
\end{itemize}

\end{enumerate}


\section{Step 2: Sum the Terms}

\begin{enumerate}

\item Write a second function that calls
\texttt{fourier\_term()} repeatedly and adds the results:

\begin{verbatim}
def sawtooth(x, N):
    total = 0.0

    for n in range(1, N + 1):
        total = total + \
                fourier_term(x, n)

    return total
\end{verbatim}

Test the function at one point:

\begin{verbatim}
x = np.pi / 4

print(sawtooth(x, 1))
print(sawtooth(x, 5))
print(sawtooth(x, 20))
\end{verbatim}

\begin{itemize}
\item Which function performs one calculation?
\item Which function controls the summation?
\item Why does \texttt{sawtooth()} call
\texttt{fourier\_term()} inside a loop?
\end{itemize}

\end{enumerate}


\section{Step 3: Plot the Series}

\begin{enumerate}

\item Create a third function whose job is only to plot the
approximation:

\begin{verbatim}
def plot_series(x, N):
    y = sawtooth(x, N)

    plt.plot(x, y,
             label=f"N = {N}")
\end{verbatim}

This function does not calculate the Fourier terms directly.
Instead, it calls \texttt{sawtooth()}, which in turn calls
\texttt{fourier\_term()}.

\end{enumerate}


\section{Step 4: Put It Together}

\begin{enumerate}

\item At the bottom of the program, write the main script:

\begin{verbatim}
x = np.linspace(-np.pi,
                np.pi, 400)

plot_series(x, 1)
plot_series(x, 5)
plot_series(x, 20)

plt.plot(x, x,
         label="f(x) = x")

plt.xlabel("x")
plt.ylabel("f(x)")
plt.title("Sawtooth Fourier Series")
plt.legend()
plt.grid()
plt.show()
\end{verbatim}

\begin{itemize}

\item Run the complete program and compare the three
approximations.

\item Which function calls \texttt{fourier\_term()}?

\item Which function calls \texttt{sawtooth()}?

\item Which part of the program controls which values of
$N$ are plotted?

\item Change the final call from

\begin{verbatim}
plot_series(x, 20)
\end{verbatim}

to

\begin{verbatim}
plot_series(x, 50)
\end{verbatim}

and observe how the approximation changes.

\end{itemize}

\end{enumerate}

\end{document}